\chapter{Les dépendances}
\label{ch:dependances}

Un projet dépend rarement de rien. Ce chapitre couvre les trois formes de
dépendance --- entre projets, vers des bibliothèques, vers des fichiers --- et la
distinction, souvent floue, entre \emph{lier} et \emph{dépendre}.

\section{Dépendre d'un autre projet}

\begin{nkcode}{}
dependson(["NKCore", "NKMath"])
\end{nkcode}

\begin{nkretenir}
\textbf{\texttt{dependson} fait deux choses à la fois}, et c'est pour cela qu'il
existe :

\begin{enumerate}
\item il impose un \textbf{ordre} --- les dépendances se construisent avant ;
\item il fait \textbf{hériter} ce que la dépendance exporte : ses chemins
      d'inclusion, ses définitions publiques.
\end{enumerate}

C'est l'héritage qui compte le plus. Sans lui, il faudrait répéter les chemins
d'inclusion de chaque bibliothèque dans chaque consommateur --- et les répéter,
c'est les laisser diverger.
\end{nkretenir}

\begin{nkretenir}
La sortie de construction montre l'ordre calculé, et c'est une ligne à lire :

\begin{nkcode}{Capture reelle}
Build Order (1 projects):
  1. Bonjour [CONSOLE_APP]
\end{nkcode}

Sur un vrai dépôt elle en énumère des dizaines. \textbf{Si un projet n'y figure
pas, il ne sera pas construit} --- quelles que soient vos modifications dans ses
sources.
\end{nkretenir}

\begin{nkpiege}
\textbf{Un cycle de dépendances n'a pas de solution.} A dépend de B qui dépend de
A : aucun ordre ne satisfait les deux.

Le remède n'est jamais de « casser » l'ordre mais de \textbf{casser le cycle} :
extraire dans un troisième projet ce dont les deux ont besoin. Un cycle est
toujours le signe qu'une chose commune n'a pas été nommée.
\end{nkpiege}

\section{Lier une bibliothèque}

\begin{nkcode}{}
links(["user32", "gdi32", "opengl32"])
libdirs(["Externals/Libs"])
\end{nkcode}

\begin{center}
\begin{tabular}{@{}ll@{}}
\toprule
\texttt{links} & les bibliothèques à lier, sans préfixe ni extension \\
\texttt{libdirs} & où les chercher \\
\texttt{syslibdirs} & idem, mais chemins système \\
\texttt{removelinks} & retire d'une liste héritée \\
\bottomrule
\end{tabular}
\end{center}

\begin{nkretenir}
\textbf{\texttt{links(["X11"])} devient \texttt{-lX11}}, et rien de plus. Jenga ne
vérifie pas que la bibliothèque existe : il passe le drapeau, et c'est l'éditeur
de liens qui échoue.

D'où la distinction qu'il faut garder en tête : \texttt{dependson} parle de
\emph{votre} dépôt, \texttt{links} parle du \emph{système}. Le premier, Jenga le
construit ; le second, il suppose qu'il est là.
\end{nkretenir}

\begin{nkretenir}[Sur macOS et iOS, ce ne sont pas des bibliothèques]
\begin{nkcode}{}
frameworks(["Cocoa", "QuartzCore", "Metal", "AudioToolbox"])
\end{nkcode}

Un \emph{framework} Apple est un dossier qui porte à la fois l'en-tête et le
binaire. Il se déclare par \texttt{frameworks}, pas par \texttt{links}.

\alerte{} \textbf{Et la liste n'est pas la même sur macOS et sur iOS.}
\texttt{AudioUnit} existe sur macOS et \textbf{pas} sur iOS, où ses symboles
vivent dans \texttt{AudioToolbox}. Recopier la liste macOS donne
\texttt{ld: framework AudioUnit not found} --- mesuré. \emph{Un précédent se
vérifie sur SA plateforme.}
\end{nkretenir}

\section{Dépendre de fichiers}

\begin{nkcode}{}
dependfiles(["assets"])
embedresources(["Resources/police.ttf"])
\end{nkcode}

\begin{center}
\begin{tabular}{@{}ll@{}}
\toprule
\texttt{dependfiles} & copiés à côté de l'artefact \\
\texttt{embedresources} & intégrés \emph{dans} le binaire \\
\bottomrule
\end{tabular}
\end{center}

\begin{nkretenir}
\textbf{La différence décide de votre portabilité.} Un fichier copié doit être
trouvé à l'exécution --- ce qui dépend du répertoire courant, des permissions
Android, et de la façon dont le Web sert ses fichiers. Un fichier intégré est
toujours là.

Les jeux de Nkentseu ne déclarent \emph{aucun} asset, et leur descripteur
explique pourquoi :

\begin{nkcode}{Applications/NkDames/NkDames.jenga}
# AUCUN dossier assets : NkDames ne charge aucun fichier a l'execution.
# Le plateau et les pieces sont dessines par geometrie, et la police est
# EMBARQUEE. C'est ce qui lui permet de demarrer a l'identique sur les six
# plateformes sans chaine de copie d'assets.
\end{nkcode}
\end{nkretenir}

\section{Les chemins d'inclusion}

\begin{center}
\begin{tabular}{@{}ll@{}}
\toprule
\texttt{includedirs} & vos en-têtes \\
\texttt{externalincludedirs} & ceux d'un tiers --- avertissements réduits \\
\texttt{sysincludedirs} & en-têtes système \\
\texttt{removeincludedirs} & retire d'une liste héritée \\
\bottomrule
\end{tabular}
\end{center}

\begin{nkretenir}
\textbf{\texttt{externalincludedirs} n'est pas un synonyme.} Il passe le chemin
comme \emph{externe} au compilateur, ce qui coupe les avertissements venant de
ces en-têtes.

Sans lui, compiler avec des avertissements exigeants noie vos propres messages
sous ceux d'une bibliothèque que vous ne corrigerez pas. \textbf{Un avertissement
qu'on ne peut pas traiter rend inutiles ceux qu'on pourrait traiter.}
\end{nkretenir}

\begin{nkpiege}
\textbf{C'est la ligne qu'on oublie, et son symptôme accuse votre code.}
Sans \texttt{extra\_includes} ou \texttt{includedirs}, vous obtenez
\texttt{fatal error: 'MonEnTete.h' file not found} --- une erreur qui pointe
\emph{votre} inclusion, alors que c'est le \emph{projet} qui ne sait pas où
chercher.
\end{nkpiege}

\section{Ce que Nkentseu ajoute, et qui n'est PAS de Jenga}

\begin{nkretenir}
Les descripteurs de Nkentseu appellent \texttt{nkentseudependson(\dots)}. Cette
fonction \textbf{n'existe pas dans Jenga} : elle est définie dans le dépôt, à
\nkfile{config/modules.jenga}, ligne 334.

\begin{nkcode}{config/modules.jenga:334}
def nkentseudependson(deps=None, selfexport=None, extra_includes=None,
                      extra_links=None, extra_defines=None) -> None:
\end{nkcode}

C'est une \textbf{démonstration de ce que le chapitre 4 annonçait} : puisqu'un
\texttt{.jenga} est du Python, un dépôt peut se donner son propre vocabulaire.
Ici, une seule fonction pose les dépendances, les chemins d'inclusion et les
définitions selon une table centrale --- au lieu de les répéter dans trente
descripteurs.

\textbf{Sachez distinguer les deux quand vous lisez un descripteur} : ce qui vient
de Jenga se retrouve dans la documentation de Jenga ; le reste est à vous, et
c'est dans votre dépôt qu'il faut le chercher.
\end{nkretenir}

\begin{nkpiege}
\textbf{Une dépendance déclarée à DEUX endroits ne se corrige pas à un seul.}
Dans Nkentseu, un module déclare ses dépendances dans son propre \texttt{.jenga}
\emph{et} dans la table centrale de \texttt{config/modules.jenga} --- et c'est la
table qui gouverne.

Mesuré : retirer une dépendance du seul \texttt{.jenga} d'un module ne l'a pas
retirée du graphe. Le défaut s'est vu en \emph{mesurant} le graphe construit, pas
en relisant les fichiers.
\end{nkpiege}

\begin{nkbilan}
Trois formes de dépendance : \texttt{dependson} entre projets --- qui impose
l'ordre \emph{et} fait hériter les inclusions ; \texttt{links} vers le système ---
que Jenga suppose présent sans le vérifier ; \texttt{dependfiles} et
\texttt{embedresources} pour les fichiers, et le choix entre les deux décide de
votre portabilité. Les frameworks Apple ne sont pas des bibliothèques, et leur
liste diffère entre macOS et iOS. Enfin, un dépôt peut se donner son propre
vocabulaire --- \texttt{nkentseudependson} n'est pas de Jenga --- et une
dépendance déclarée à deux endroits ne se corrige pas à un seul.
\end{nkbilan}

\begin{nkquestions}
\begin{enumerate}
\item Quelles \emph{deux} choses \texttt{dependson} fait-il, et laquelle compte
      le plus ?
\item Quelle différence entre \texttt{dependson} et \texttt{links} ?
\item Que devient \texttt{links(["X11"])}, et qui échoue si la bibliothèque
      manque ?
\item Pourquoi \texttt{externalincludedirs} plutôt que \texttt{includedirs} pour
      un tiers ?
\item \texttt{nkentseudependson} est-il une fonction de Jenga ? Où la trouve-t-on ?
\end{enumerate}
\end{nkquestions}

\begin{nkpratique}
Créez trois projets : une bibliothèque statique \texttt{Calcul}, une seconde
\texttt{Affichage} qui en dépend, et un exécutable qui dépend des deux.

\begin{enumerate}
\item Vérifiez l'ordre dans la ligne \texttt{Build Order}.
\item Placez un en-tête dans \texttt{Calcul} et incluez-le depuis l'exécutable
      \emph{sans} ajouter de \texttt{includedirs} --- que se passe-t-il, et
      pourquoi ?
\item Créez volontairement un cycle et lisez ce que Jenga dit.
\end{enumerate}
\end{nkpratique}

\begin{nkdemo}
Prouvez que \texttt{dependson} fait hériter, et pas seulement ordonner.

Deux projets : une bibliothèque qui expose un en-tête, un exécutable qui
l'inclut. Faites-le fonctionner avec \texttt{dependson}, puis \textbf{retirez le
\texttt{dependson} en gardant \texttt{links}} sur la bibliothèque.

\textbf{Ce que la démonstration doit établir} : l'échec change de \emph{nature}.
Avec l'héritage, tout passe ; sans lui, l'erreur n'est pas à l'édition de liens
--- où la bibliothèque est pourtant toujours donnée --- mais à la
\emph{compilation}, sur l'en-tête introuvable.

Relevez les deux messages. C'est leur différence qui prouve que l'héritage porte
les chemins d'inclusion, pas les symboles.
\end{nkdemo}

\begin{nklogique}
Vous retirez une dépendance du \texttt{.jenga} d'un module, vous reconstruisez, et
le graphe la contient toujours.

\begin{enumerate}
\item Énumérez trois explications possibles.
\item Pour chacune, dites quelle commande la confirmerait --- et laquelle ne se
      voit \emph{pas} en relisant le fichier que vous venez de modifier.
\item Concluez sur la règle générale : que faut-il faire avant de croire qu'une
      suppression a pris effet ?
\end{enumerate}
\end{nklogique}
